\documentclass{stex}

\libusepackage{naproche}
\libinput{preamble}

\begin{document}
\begin{smodule}{modular-arithmetics.ftl}

\importmodule[libraries/arithmetics]{multiplication-and-ordering.ftl}

\symdef{Div}{\,\textrm{div}\,}
\symdef{Mod}{\,\textrm{mod}\,}
\symdef{Cong}[args=3]{#1\;\comp\equiv\;#2\;\comp{(\textrm{mod}}\,#3\comp)}

\symdecl*{quotient}
\symdecl*{remainder}
\symdecl*{congruent}
\symdecl*{congruency of sum}

\begin{definition}[forthel,for={Div,quotient}]
  Let $n, m$ be natural numbers such that $m \NotEq \Zero$.
  $\emph{n \Div m}$ is the natural number $q$ such that $n \Eq (m \Mul q) \Plus r$ for some natural number $r$ that is less than $m$.
  Let the \emph{quotient of $n$ over $m$} stand for $n \Div m$.
\end{definition}

\begin{definition}[forthel,for={Mod,remainder}]
  Let $n, m$ be natural numbers such that $m \NotEq \Zero$.
  $\emph{n \Mod m}$ is the natural number $r$ such that $r \Less m$ and there exists a natural number $q$ such that $n \Eq (m \Mul q) \Plus r$.
  Let the \emph{remainder of $n$ over $m$} stand for $n \Mod m$.
\end{definition}

\begin{definition}[forthel,for={Cong,congruent}]
  Let $n, m, k$ be natural numbers such that $k \NotEq \Zero$.
  $\emph{\Cong{n}{m}{k}}$ iff $n \Mod k \Eq m \Mod k$.
  Let $n$ and $m$ are \emph{congruent modulo $k$} stand for $\Cong{n}{m}{k}$.
\end{definition}

\begin{proposition}[forthel]
  Let $n, k$ be natural numbers such that $k \NotEq \Zero$.
  Then $\Cong{n}{n}{k}$.
\end{proposition}
\begin{proof}[forthel]
  We have $n \Mod k \Eq n \Mod k$.
  Hence $\Cong{n}{n}{k}$.
\end{proof}

\begin{proposition}[forthel]
  Let $n, m, k$ be natural numbers such that $k \NotEq \Zero$.
  If $\Cong{n}{m}{k}$ then $\Cong{m}{n}{k}$.
\end{proposition}
\begin{proof}[forthel]
  Assume $\Cong{n}{m}{k}$.
  Then $n \Mod k \Eq m \Mod k$.
  Hence $m \Mod k \Eq n \Mod k$.
  Thus $\Cong{m}{n}{k}$.
\end{proof}

\begin{proposition}[forthel]
  Let $n, m, k$ be natural numbers such that $k \NotEq \Zero$.
  If $\Cong{n}{m}{k}$ and $\Cong{m}{\One}{k}$ then $\Cong{n}{\One}{k}$.
\end{proposition}
\begin{proof}[forthel]
  Assume $\Cong{n}{m}{k}$ and $\Cong{m}{\One}{k}$.
  Then $n \Mod k \Eq m \Mod k$ and $m \Mod k \Eq \One \Mod k$.
  Hence $n \Mod k \Eq \One \Mod k$.
  Thus $\Cong{n}{\One}{k}$.
\end{proof}

\begin{proposition}[forthel]
  Let $n, m, k$ be natural numbers such that $k \NotEq \Zero$.
  Assume $n \Geq m$.
  Then $\Cong{n}{m}{k}$ iff $n \Eq (k \Mul x) \Plus m$ for some natural number $x$.
\end{proposition}
\begin{proof}[forthel]
  \begin{case}{$\Cong{n}{m}{k}$.}
    Then $n \Mod k \Eq m \Mod k$.
    Take a natural number $r$ such that $r \Less k$ and $n \Mod k \Eq r \Eq m \Mod k$.
    Take a nonzero natural number $l$ such that $k \Eq r \Plus l$.
    Consider natural numbers $q,q'$ such that $n \Eq (q \Mul k) \Plus r$ and $m \Eq (q' \Mul k) \Plus r$.

    Then $q \Geq q'$.
    \begin{proof}
      Assume the contrary.
      Then $q \Less q'$.
      Hence $q \Mul k \Less q' \Mul k$.
      Thus $(q \Mul k) \Plus r \Less (q' \Mul k) \Plus r$ (by \sn{preservation of ordering under right-addition}).
      Indeed $q \Mul k$ and $q' \Mul k$ are natural numbers.
      Therefore $n \Less m$.
      Contradiction.
    \end{proof}

    Take a natural number $x$ such that $q \Eq q' \Plus x$.

    Let us show that $n \Eq (k \Mul x) \Plus m$.
      We have
      \[  (k \Mul x) \Plus m                       \]
      \[    \Eq (k \Mul x) \Plus ((q' \Mul k) \Plus r)  \]
      \[    \Eq ((k \Mul x) \Plus (q' \Mul k)) \Plus r  \]
      \[    \Eq ((k \Mul x) \Plus (k \Mul q')) \Plus r  \]
      \[    \Eq (k \Mul (q' \Plus x)) \Plus r            \]
      \[    \Eq (k \Mul q) \Plus r                   \]
      \[    \Eq n.                                \]
    End.
  \end{case}

  \begin{case}{$n \Eq (k \Mul x) \Plus m$ for some natural number $x$.}
    Consider a natural number $x$ such that $n \Eq (k \Mul x) \Plus m$.
    Take natural numbers $r, r'$ such that $n \Mod k \Eq r$ and $m \Mod k \Eq r'$.
    Then $r, r' \Less k$.
    Take natural numbers $q, q'$ such that $n \Eq (k \Mul q) \Plus r$ and $m \Eq (k \Mul q') \Plus r'$.
    Then
    \[  (k \Mul q) \Plus r                         \]
    \[    \Eq n                                   \]
    \[    \Eq (k \Mul x) \Plus m                     \]
    \[    \Eq (k \Mul x) \Plus ((k \Mul q') \Plus r')   \]
    \[    \Eq ((k \Mul x) \Plus (k \Mul q')) \Plus r'   \]
    \[    \Eq (k \Mul (x \Plus q')) \Plus r'.            \]
    Hence $r \Eq r'$.
    Thus $n \Mod k \Eq m \Mod k$.
    Therefore $\Cong{n}{m}{k}$.
  \end{case}
\end{proof}

\begin{proposition}[forthel]
  Let $n, m, k, k'$ be natural numbers such that $k, k' \NotEq \Zero$.
  If $\Cong{n}{m}{k \Mul k'}$ then $\Cong{n}{m}{k}$.
\end{proposition}
\begin{proof}[forthel]
  Assume $\Cong{n}{m}{k \Mul k'}$.

  \begin{case}{$n \Geq m$.}
    We can take a natural number $x$ such that $n \Eq ((k \Mul k') \Mul x) \Plus m$.
    Then $n \Eq (k \Mul (k' \Mul x)) \Plus m$.
    Hence $\Cong{n}{m}{k}$.
  \end{case}

  \begin{case}{$m \Geq n$.}
    We have $\Cong{m}{n}{k \Mul k'}$.
    Hence we can take a natural number $x$ such that $m \Eq ((k \Mul k') \Mul x) \Plus n$.
    Then $m \Eq (k \Mul (k' \Mul x)) \Plus n$.
    Thus $\Cong{m}{n}{k}$.
    Therefore $\Cong{n}{m}{k}$.
  \end{case}
\end{proof}

\begin{corollary}[forthel]
  Let $n, m, k, k'$ be natural numbers such that $k, k' \NotEq \Zero$.
  If $\Cong{n}{m}{k \Mul k'}$ then $\Cong{n}{m}{k'}$.
\end{corollary}
\begin{proof}[forthel]
  Assume $\Cong{n}{m}{k \Mul k'}$.
  Then $\Cong{n}{m}{k' \Mul k}$.
  Hence $\Cong{n}{m}{k'}$.
\end{proof}

\begin{proposition}[forthel,name=congruency of sum]
  Let $n, k$ be natural numbers such that $k \NotEq \Zero$.
  Then $\Cong{n \Plus k}{n}{k}$.
\end{proposition}
\begin{proof}[forthel]
  Take $r \Eq n \Mod k$ and $r' \Eq (n \Plus k) \Mod k$.
  Consider a $q \Elem \Naturals$ such that $n \Eq (k \Mul q) \Plus r$ and $r \Less k$.
  Consider a $q' \Elem \Naturals$ such that $n \Plus k \Eq (k \Mul q') \Plus r'$ and
  $r' \Less k$.
  Then $(k \Mul q') \Plus r'
    \Eq n \Plus k
    \Eq ((k \Mul q) \Plus r) \Plus k
    \Eq (k \Plus (k \Mul q)) \Plus r
    \Eq (k \Mul (q \Plus \One)) \Plus r$.
  Hence $r \Eq r'$.
  Consequently $n \Mod k \Eq (n \Plus k) \Mod k$.
  Thus $\Cong{n \Plus k}{n}{k}$.
\end{proof}

\end{smodule}
\end{document}
